% Part: first-order-logic
% Chapter: syntax-and-semantics
% Section: formation-sequences

\documentclass[../../../include/open-logic-section]{subfiles}

\begin{document}

\olfileid[pt]{fol}{syn}{fseq}

\olsection{Sequências de Formação}

Definir !!{formula}s via uma definição indutiva, e a
técnica complementar de provar propriedades de !!{formula}s via
indução, é uma abordagem elegante e eficiente. No entanto, também pode
ser útil considerar uma abordagem mais ascendente, passo a passo,
da construção de !!{formula}s, o que fazemos aqui usando a
noção de uma \emph{sequência de formação}.
%
Para mostrar como termos e !!{formula}s podem ser introduzidos dessa maneira
sem precisar referir-se a suas definições indutivas, primeiro
introduzimos a noção de uma cadeia arbitrária de símbolos extraída de
alguma linguagem~$\Lang L$.

\begin{defn}[Cadeias]
\ollabel{defn:string}
Suponha que $\Lang L$ é uma linguagem de primeira ordem. Uma \emph{$\Lang
L$-cadeia} é uma sequência finita de símbolos de~$\Lang L$. Quando a
linguagem~$\Lang L$ está claramente fixada pelo contexto, frequentemente nos
referiremos a uma $\Lang L$-cadeia simplesmente como uma \emph{cadeia}.
\end{defn}

\begin{ex}
Para qualquer linguagem de primeira ordem $\Lang L$, todas as
$\Lang L$-!!{formula}s são $\Lang L$-cadeias, mas não
reciprocamente. Por exemplo, \[)(\Obj v_0\lif\lexists\] é uma
$\Lang L$-cadeia mas não uma $\Lang L$-!!{formula}.
\end{ex}

\begin{defn}[Sequências de formação para termos]
\ollabel{defn:fseq-trm}
Uma sequência finita de $\Lang L$-cadeias $\tuple{t_0,\dotsc,t_n}$ é uma
\emph{sequência de formação} para um termo $t$ se $t \ident t_n$ e para todo
$i \leq n$, ou $t_i$ é !!a{variable} ou !!a{constant}, ou $\Lang
L$ contém um !!{function}~$f$ $k$-ário e existem
$m_0,\dotsc,m_k < i$ tais que $t_i \ident f(t_{m_0},\dotsc,t_{m_k})$.
Quando for necessário distinguir, nos referiremos às sequências de formação
para termos como \emph{sequências de formação de termos}.
\end{defn}

\begin{ex}
A sequência
\[
    \tuple{\Obj c_0, \Obj v_0, \Atom{\Obj f^2_0}{\Obj c_0, \Obj v_0}, \Atom{\Obj f^1_0}{\Atom{\Obj f^2_0}{\Obj c_0, \Obj v_0}}}
\]
é uma sequência de formação para o termo $\Atom{\Obj f^1_0}{\Atom{\Obj
f^2_0}{\Obj c_0, \Obj v_0}}$, assim como
\[
    \tuple{\Obj v_0, \Obj c_0, \Atom{\Obj f^2_0}{\Obj c_0, \Obj v_0}, \Atom{\Obj f^1_0}{\Atom{\Obj f^2_0}{\Obj c_0, \Obj v_0}}}.
\]
\end{ex}

\begin{defn}[Sequências de formação para fórmulas]
\ollabel{defn:fseq-frm}
Uma sequência finita de $\Lang L$-cadeias $\tuple{!A_0,\dotsc,!A_n}$
é uma \emph{sequência de formação} para~$!A$ se $!A \ident !A_n$ e
para todo $i \leq n$, ou $!A_i$ é uma !!{formula} atômica ou
existem $j,k < i$ e !!a{variable}~$x$ tais que uma das seguintes condições
vale:
\begin{enumerate}
    \tagitem{prvNot}{$!A_i \ident \lnot !A_j$.}{}%
    \tagitem{prvAnd}{$!A_i \ident (!A_j \land !A_k)$.}{}%
    \tagitem{prvOr}{$!A_i \ident (!A_j \lor !A_k)$.}{}%
    \tagitem{prvIf}{$!A_i \ident (!A_j \lif !A_k)$.}{}%
    \tagitem{prvIff}{$!A_i \ident (!A_j \liff !A_k)$.}{}%
    \tagitem{prvAll}{$!A_i \ident \lforall[x][!A_j]$.}{}%
    \tagitem{prvEx}{$!A_i \ident \lexists[x][!A_j]$.}{}%
\end{enumerate}
Quando for necessário distinguir, nos referiremos às sequências de formação
para fórmulas como \emph{sequências de formação de fórmulas}.
\end{defn}

\begin{ex}
\[
\tuple{
    \Atom{\Obj A^1_0}{\Obj v_0},
    \Atom{\Obj A^1_1}{\Obj c_1},
    (\Atom{\Obj A^1_1}{\Obj c_1} \land \Atom{\Obj A^1_0}{\Obj v_0}),
    \lexists[\Obj v_0][(\Atom{\Obj A^1_1}{\Obj c_1} \land \Atom{\Obj A^1_0}{\Obj v_0})]
}
\]
é uma sequência de formação de $\lexists[\Obj v_0][(\Atom{\Obj A^1_1}{\Obj
c_1} \land \Atom{\Obj A^1_0}{\Obj v_0})]$, assim como
\begin{multline*}
\tuple{
    \Atom{\Obj A^1_0}{\Obj v_0},
    \Atom{\Obj A^1_1}{\Obj c_1},
    (\Atom{\Obj A^1_1}{\Obj c_1} \land \Atom{\Obj A^1_0}{\Obj v_0}),
    \Atom{\Obj A^1_1}{\Obj c_1},\\
    \lforall[\Obj v_1][\Atom{\Obj A^1_0}{\Obj v_0}],
    \lexists[\Obj v_0][(\Atom{\Obj A^1_1}{\Obj c_1} \land \Atom{\Obj A^1_0}{\Obj v_0})]
}.
\end{multline*}
%
Como pode ser visto no segundo exemplo, sequências de formação
podem conter ``lixo'': !!{formula}s que são redundantes ou não
contribuem para a construção.
\end{ex}

\begin{prop}\ollabel{prop:formed}
Toda !!{formula}~$!A$ em~$\Frm[L]$ tem uma sequência de formação.
\end{prop}

\begin{proof}
Suponha que $!A$ é atômica. Então a sequência $\tuple{!A}$ é uma
sequência de formação para~$!A$.
%
Agora suponha que $!B$ e~$!C$ têm sequências de formação
$\tuple{!B_0,\dotsc,!B_n}$ e $\tuple{!C_0,\dotsc,!C_m}$
respectivamente.
%
\begin{enumerate}
  \tagitem{prvNot}{Se $!A \ident \lnot !B$,
      então $\tuple{!B_0,\dotsc,!B_n,\lnot !B_n}$
      é uma sequência de formação para~$!A$.}{}
  \tagitem{prvAnd}{Se $!A \ident (!B \land !C)$,
      então $\tuple{!B_0,\dotsc,!B_n,!C_0,\dotsc,!C_m,(!B_n \land !C_m)}$
      é uma sequência de formação para~$!A$.}{}
  \tagitem{prvOr}{Se $!A \ident (!B \lor !C)$,
      então $\tuple{!B_0,\dotsc,!B_n,!C_0,\dotsc,!C_m,(!B_n \lor !C_m)}$
      é uma sequência de formação para~$!A$.}{}
  \tagitem{prvIf}{Se $!A \ident (!B \lif !C)$,
      então $\tuple{!B_0,\dotsc,!B_n,!C_0,\dotsc,!C_m,(!B_n \lif !C_m)}$
      é uma sequência de formação para~$!A$.}{}
  \tagitem{prvIff}{Se $!A \ident (!B \liff !C)$,
      então $\tuple{!B_0,\dotsc,!B_n,!C_0,\dotsc,!C_m,(!B_n \liff !C_m)}$
      é uma sequência de formação para~$!A$.}{}
  \tagitem{prvAll}{Se $!A \ident \lforall[x][!B]$,
      então $\tuple{!B_0,\dotsc,!B_n,\lforall[x][!B_n]}$
      é uma sequência de formação para~$!A$.}{}
  \tagitem{prvEx}{Se $!A \ident \lexists[x][!B]$,
      então $\tuple{!B_0,\dotsc,!B_n,\lexists[x][!B_n]}$
      é uma sequência de formação para~$!A$.}{}
\end{enumerate}
Pelo princípio de indução sobre !!{formula}s,
toda !!{formula} tem uma sequência de formação.
\end{proof}

Também podemos provar a recíproca. Isso é importante porque mostra
que nossas duas maneiras de definir fórmulas são equivalentes: elas dão
os mesmos resultados. Isso também significa que podemos provar teoremas sobre
fórmulas usando indução ordinária sobre o comprimento das sequências de
formação.

\begin{lem}
\ollabel{lem:fseq-init}
Suponha que $\tuple{!A_0,\dotsc,!A_n}$ é uma sequência de formação
para~$!A_n$, e que $k \leq n$. Então $\tuple{!A_0,\dotsc,!A_k}$
é uma sequência de formação para~$!A_k$.
\end{lem}

\begin{proof}
Exercício.
\end{proof}

\begin{prob}
Prove \olref[fol][syn][fseq]{lem:fseq-init}.
\end{prob}

\begin{thm}
\ollabel{thm:fseq-frm-equiv}
$\Frm[L]$ é o conjunto de todas as $\Lang L$-cadeias~$!A$ tais que
existe uma sequência de formação de fórmulas para~$!A$.
\end{thm}

\begin{proof}
Seja $F$ o conjunto de todas as cadeias de símbolos na linguagem~$\Lang L$
que têm uma sequência de formação. Vimos em
\olref[fol][syn][fseq]{prop:formed} que $\Frm[L] \subseteq F$, então agora
provamos a recíproca.

Suponha que $!A$ tem uma sequência de formação $\tuple{!A_0,\dotsc,!A_n}$.
Provamos que $!A \in \Frm[L]$ por indução forte sobre~$n$.
Nossa hipótese de indução é que toda cadeia de símbolos com uma
sequência de formação de comprimento $m < n$ está em $\Frm[L]$.
Pela definição de uma sequência de formação, ou $!A \ident !A_n$ é
atômica ou devem existir $j,k < n$ tais que um dos
seguintes casos vale:
\begin{enumerate}
    \tagitem{prvNot}{$!A \ident \lnot !A_j$.}{}%
    \tagitem{prvAnd}{$!A \ident (!A_j \land !A_k)$.}{}%
    \tagitem{prvOr}{$!A \ident (!A_j \lor !A_k)$.}{}%
    \tagitem{prvIf}{$!A \ident (!A_j \lif !A_k)$.}{}%
    \tagitem{prvIff}{$!A \ident (!A_j \liff !A_k)$.}{}%
    \tagitem{prvAll}{$!A \ident \lforall[x][!A_j]$.}{}%
    \tagitem{prvEx}{$!A \ident \lexists[x][!A_j]$.}{}%
\end{enumerate}
Agora raciocinamos por casos. Se $!A$ é atômica então
$!A_n \in \Frm[L_0]$. Suponha, em vez disso, que
$!A \equiv (!A_j \land !A_k)$. Por 
\olref[fol][syn][fseq]{lem:fseq-init},
$\tuple{!A_0,\dotsc,!A_j}$ e $\tuple{!A_0,\dotsc,!A_k}$ são
sequências de formação para $!A_j$ e~$!A_k$, respectivamente. Como
estas são subsequências iniciais próprias da sequência de formação
para~$!A$, ambas têm comprimento menor que~$n$. Portanto, pela
hipótese de indução, $!A_j$ e~$!A_k$ estão em~$\Frm[L_0]$,
e pela definição de !!a{formula}, $(!A_j \land !A_k)$ também
está. Os outros casos seguem por raciocínio análogo.
\end{proof}

Sequências de formação para termos têm propriedades semelhantes às
das !!{formula}s.

\begin{prop}
\ollabel{prop:fseq-trm-equiv}
$\Trm[L]$ é o conjunto de todas as $\Lang L$-cadeias $t$
tais que existe uma sequência de formação de termos para~$t$.
\end{prop}

\begin{proof}
Exercício.
\end{proof}

\begin{prob}
Prove \olref[fol][syn][fseq]{prop:fseq-trm-equiv}.
Dica: use uma estratégia semelhante à usada na prova de
\olref[fol][syn][fseq]{thm:fseq-frm-equiv}.
\end{prob}

Há dois tipos de ``lixo'' que podem aparecer em sequências de
formação: elementos repetidos, e elementos que são irrelevantes
para a construção da fórmula ou termo. Podemos eliminar
ambos olhando para sequências de formação mínimas.

\begin{defn}[Sequências de formação mínimas]
\ollabel{defn:minimal-fseq}
Uma sequência de formação $\tuple{!A_0, \ldots, !A_n}$ para uma
fórmula~$!A$ é uma \emph{sequência de formação mínima} para~$!A$
se, para toda outra sequência de formação~$s$ para~$!A$,
o comprimento de~$s$ é maior ou igual a~$n+1$.

Similarmente, uma sequência de formação $\tuple{t_0, \ldots, t_n}$
para um termo~$t$ é uma \emph{sequência de formação mínima}
para~$t$ se, para toda outra sequência de formação~$s$ para~$t$,
o comprimento de~$s$ é maior ou igual a~$n+1$.
\end{defn}

Note que uma fórmula ou termo pode ter mais de uma sequência de formação
mínima, mas elas conterão exatamente as mesmas
cadeias.

\begin{prop}
\ollabel{prop:subformula-equivs}
As seguintes condições são equivalentes:
\begin{enumerate}
    \item $!B$ é uma sub-!!{formula} de~$!A$.
    \item $!B$ ocorre em toda sequência de formação de~$!A$.
    \item $!B$ ocorre em uma sequência de formação mínima de~$!A$.
\end{enumerate}
\end{prop}

\begin{proof}
Exercício.
\end{proof}

\begin{prob}
Prove \olref[fol][syn][fseq]{prop:subformula-equivs}.
\end{prob}

\begin{history}
Sequências de formação foram introduzidas por Raymond Smullyan em seu
livro-texto \emph{First-Order Logic} \citep{Smullyan1968}.
Propriedades adicionais de sequências de formação foram estabelecidas por
\citet{Zuckerman1973}.
\end{history}

\end{document}
